% ============================================================================
% LANGENTWURF LE-011: Browser & Internet
% Profil: S (Senioren 65+) | Tier: 2 (Standard)
% ============================================================================

\documentclass[11pt,a4paper]{article}
\usepackage[utf8]{inputenc}
\usepackage[T1]{fontenc}
\usepackage[ngerman]{babel}
\usepackage{geometry}
\usepackage{fancyhdr}
\usepackage{lastpage}
\usepackage{xcolor}
\usepackage{tcolorbox}
\usepackage{enumitem}
\usepackage{longtable}
\usepackage{hyperref}
\usepackage{amssymb}

\geometry{left=2.5cm, right=2.5cm, top=2.5cm, bottom=2.5cm}
\definecolor{fach}{HTML}{b35c00}
\definecolor{gray}{HTML}{666666}
\definecolor{successgreen}{HTML}{2a7a4b}

\tcbuselibrary{skins,breakable}
\newtcolorbox{scorebox}{ colback=white, colframe=fach, fonttitle=\bfseries, title=Didaktik-Score, breakable }

\pagestyle{fancy}
\fancyhf{}
\fancyhead[L]{\small\textcolor{gray}{Digitale Grundlagen -- 011 Browser \& Internet}}
\fancyhead[R]{\small\textcolor{gray}{LE Tier 2 / Profil S}}
\fancyfoot[C]{\small\thepage{} / \pageref{LastPage}}
\renewcommand{\headrulewidth}{0.4pt}

\begin{document}

\begin{titlepage}
    \centering
    \vspace*{2cm}
    {\Large\textcolor{gray}{Digitale Grundlagen -- Senioren 65+}}\\[2cm]
    {\Huge\textcolor{fach}{\textbf{Langentwurf}}}\\[0.5cm]
    {\large Tier 2 | Profil S}\\[2cm]
    {\LARGE\textbf{011 -- Browser \& Internet}}\\[0.5cm]
    {\large Modul C: Internet \& Kommunikation}\\[2cm]
    \begin{tabular}{rl}
        \textbf{Zeitrahmen:} & 60--75 Minuten \\
        \textbf{Vorkenntnisse:} & WLAN verbinden, Apps öffnen \\
    \end{tabular}
    \vfill
    {\normalsize Stand: \today}
\end{titlepage}

\tableofcontents
\newpage

\section{Bedingungsanalyse}

\subsection{Ausgangssituation}
\begin{itemize}
    \item Safari ist oft unbekannt oder wird mit ``Internet'' gleichgesetzt
    \item Tabs werden nicht verstanden (viele offene Seiten)
    \item Lesezeichen werden nicht genutzt
    \item Adressleiste vs. Suchfeld ist unklar
    \item Pop-ups und Werbung verwirren
\end{itemize}

\subsection{Typische Probleme}
\begin{itemize}
    \item ``Wie komme ich auf eine Webseite?''
    \item ``Wo ist die Seite hin, die ich gerade hatte?''
    \item ``Mein iPhone ist langsam'' (100 offene Tabs)
    \item ``Wie speichere ich eine Seite?''
\end{itemize}

\section{Sachanalyse}

\subsection{Kernkonzepte}
\begin{enumerate}
    \item \textbf{Browser (Safari):} Das ``Fenster ins Internet'' -- Programm zum Anzeigen von Webseiten
    \item \textbf{Webadresse (URL):} Die ``Adresse'' einer Webseite (z.B. www.google.de)
    \item \textbf{Suchmaschine:} Findet Webseiten für dich (Google, Bing)
    \item \textbf{Tabs:} Mehrere Seiten gleichzeitig offen (wie Karteikarten)
    \item \textbf{Lesezeichen:} Gespeicherte Favoriten-Seiten
    \item \textbf{Verlauf:} Liste der besuchten Seiten
\end{enumerate}

\subsection{Alltagsanalogien}
\begin{tabular}{|l|p{8cm}|}
\hline
\textbf{Begriff} & \textbf{Analogie} \\
\hline
Browser & Fenster, durch das du die Welt (Internet) siehst \\
Webadresse & Straßenadresse einer Firma \\
Suchmaschine & Telefonbuch, das Adressen für dich sucht \\
Tabs & Mehrere Bücher gleichzeitig aufgeschlagen \\
Lesezeichen & Lesezeichen im Buch -- markiert wichtige Stellen \\
Verlauf & Logbuch deiner Reisen \\
\hline
\end{tabular}

\section{Didaktische Analyse}

\subsection{Stellung in der Reihe}
Erste Einheit in Modul C (Internet \& Kommunikation). Grundlage für Suchen, E-Mail, Online-Banking.

\subsection{Didaktische Reduktion}
\textbf{Ausgeklammert:} Private Browsing, Downloads, Reader-Ansicht, Browser-Erweiterungen.

\textbf{Fokus:} Grundlegende Navigation im Internet mit Safari.

\section{Lernziele}

\begin{tabular}{|c|p{7cm}|c|c|}
\hline
\textbf{Nr.} & \textbf{Die Teilnehmer können...} & \textbf{Stufe} & \textbf{Niveau} \\
\hline
FZ1 & Safari öffnen und eine Webadresse eingeben & Anwenden & Basis \\
FZ2 & über die Suchleiste nach Begriffen suchen & Anwenden & Basis \\
FZ3 & zwischen mehreren Tabs wechseln und Tabs schließen & Anwenden & Basis \\
FZ4 & eine Webseite als Lesezeichen speichern & Anwenden & Basis \\
FZ5 & den Verlauf aufrufen und eine frühere Seite finden & Anwenden & Vertiefung \\
\hline
\end{tabular}

\section{Verlaufsplanung}

\begin{longtable}{|p{1cm}|p{2cm}|p{5cm}|p{3cm}|p{2cm}|}
\hline
\textbf{Zeit} & \textbf{Phase} & \textbf{Inhalt} & \textbf{TN-Aktivität} & \textbf{Material} \\
\hline
0--10 & Einstieg & ``Was ist das Internet?'' Begriffe sammeln, Safari zeigen & Nachdenken & -- \\
\hline
10--25 & Webseite öffnen & Adressleiste, Webadresse eingeben (z.B. www.tagesschau.de) & Mitmachen & S.1 \\
\hline
25--35 & Suchen & Google-Suche über Adressleiste, Ergebnisse verstehen & Selbst suchen & S.1 \\
\hline
35--50 & Tabs & Tabs erklären, neuer Tab, wechseln, schließen, alle schließen & Üben & S.2 \\
\hline
50--60 & Lesezeichen & Lesezeichen hinzufügen, aufrufen, organisieren & Mitmachen & S.2 \\
\hline
60--70 & Übungen & Checkliste durcharbeiten, Verlauf zeigen & Bearbeiten & S.3 \\
\hline
\end{longtable}

\section{Lösungen}

\textbf{Safari-Bedienung:}
\begin{itemize}
    \item Adressleiste: Oben, dort Adresse ODER Suchbegriff eingeben
    \item Neuer Tab: + Symbol antippen
    \item Tabs anzeigen: Zwei überlappende Quadrate antippen
    \item Tab schließen: X am Tab oder zur Seite wischen
    \item Lesezeichen hinzufügen: Teilen-Symbol (Quadrat mit Pfeil) $\rightarrow$ Lesezeichen
\end{itemize}

\textbf{Wichtige Webseiten für Senioren:}
\begin{itemize}
    \item www.tagesschau.de -- Nachrichten
    \item www.wetter.de -- Wetter
    \item www.apotheken-umschau.de -- Gesundheit
    \item www.bahn.de -- Fahrpläne
\end{itemize}

\section{Didaktik-Score}

\begin{scorebox}
\begin{tabular}{|c|l|c|c|p{3.5cm}|}
\hline
\# & Dimension & Gew. & Score & Notiz \\
\hline
1 & Differenzierung & 15\% & 8/10 & Basis/Vertiefung klar \\
2 & Taxonomie & 10\% & 9/10 & Anwenden dominiert \\
3 & Alltagsrelevanz & 10\% & 10/10 & Internet-Zugang fundamental \\
4 & Aktivierung & 10\% & 10/10 & Ständiges Mitmachen \\
5 & Scaffolding & 10\% & 9/10 & Gute Analogien \\
6 & Feedback & 10\% & 8/10 & Checkliste vorhanden \\
7 & Sprachliche Klarheit & 10\% & 9/10 & Einfache Erklärungen \\
8 & Motivation & 10\% & 9/10 & Praktische Seiten zeigen \\
9 & Barrierefreiheit & 10\% & 9/10 & Große Darstellung \\
10 & Praktikabilität & 5\% & 9/10 & 70 Min angemessen \\
\hline
\multicolumn{3}{|r|}{\textbf{GESAMT:}} & \textbf{9.05/10} & \textcolor{successgreen}{\textbf{FREIGABE}} \\
\hline
\end{tabular}
\end{scorebox}

\end{document}
